% ============================================================================
% Fichier  : partie6/P06_C28_analyse_primitives.tex
% Titre    : Analyse des Primitives Cryptographiques selon le Cadre CLO
% Auteur   : MINKA MI NGUIDJOÏ Thierry Emmanuel
% Version  : 2.0 -- Juin 2026 
% ============================================================================

\chapter{Analyse des Primitives Cryptographiques selon le Cadre CLO}
\label{chap:prim-crypt}

\epigraph{%
    « Cryptographic primitives are the building blocks of trust in
    digital systems. Each must be evaluated against its resilience to
    emerging threats. »%
}{Adi Shamir}

\niveauChapitre{Recherche}{%
    Chapitre~\ref{chap:CRO} (formalisation du Trilemme CRO)
    indispensable~: ce chapitre applique méthodiquement les outils
    qui y sont définis. Notions de cryptographie avancée recommandées.%
}

\begin{boxobjectifs}
\begin{itemize}
    \item Reproduire, étape par étape, la méthodologie de calcul du
          cadre CLO sur les trois protocoles effectivement analysés
          dans la littérature publiée~: Helios, zk-SNARKs, ECDSA.
    \item Comprendre les trois méthodes d'estimation de l'entropie
          contextuelle $H(C)$~: bornes structurelles, littérature
          d'expertise, modélisation information-théorique.
    \item Décomposer le calcul de $I(V;\Sigma\mid C)$ dimension par
          dimension pour un protocole donné.
    \item Situer chaque primitive dans son \textbf{régime de capacité
          institutionnelle} (faible/moyen/élevé) et en déduire des
          recommandations de conception.
    \item Distinguer rigoureusement ce qui est \textbf{publié et
          calculé} de ce qui reste, en l'état de la recherche,
          \textbf{une question ouverte} pour d'autres primitives.
    \item Appliquer cette méthodologie à une primitive non encore
          analysée dans la littérature, comme exercice de recherche.
\end{itemize}
\end{boxobjectifs}

\bigskip

% ============================================================================
\section*{De la théorie générale à l'analyse de cas reproductible}
% ============================================================================

Le Chapitre~\ref{chap:CRO} a établi le cadre théorique~: définitions
formelles, théorèmes de bornes, démonstration du Trilemme conditionnel.
Ce chapitre descend d'un niveau d'abstraction~: il \textbf{reproduit
intégralement} la méthodologie de calcul appliquée par les auteurs
à trois protocoles réels, avec une transparence totale sur chaque
étape, l'objectif pédagogique étant de vous rendre capable de
mener vous-même ce type d'analyse sur une primitive de votre choix.

\begin{boxalert}
\textbf{Avertissement méthodologique, conforme à celui des auteurs}

Les calculs présentés dans ce chapitre sont qualifiés, dans
l'article source~\cite{minka2026clo}, d'\textbf{estimations
théoriques} (\textit{theoretical estimates}), pas de mesures
empiriques. La note méthodologique de l'article le précise
explicitement~: \textit{``ces estimations démontrent une capacité
analytique et identifient des motifs de compromis~; la validation
empirique auprès de décideurs institutionnels reste un travail
futur''}. Ce chapitre conserve scrupuleusement cette honnêteté
méthodologique, exactement la posture de la distinction
constatation/analyse pratiquée depuis le Ch.~\ref{chap:rapport}
de ce cours.
\end{boxalert}

\separateur

% ============================================================================
\section{Méthodologie de Calcul~: les Quatre Étapes}
\label{sec:methodologie}
% ============================================================================

Pour chaque protocole analysé, la méthodologie publiée
dans~\cite{minka2026clo} procède en quatre étapes systématiques.

\begin{tableaucours}{p{2.5cm}p{8.5cm}}{%
    \tableauHeader{Étape} &
    \tableauHeader{Contenu}
}
\textbf{1. Formalisation} & Définir $\mathcal{M}, \Sigma, \mathcal{C},
    \mathcal{V}$ à partir des spécifications publiées du protocole. \\
\textbf{2. Estimation de $H(C)$} & Estimer l'entropie contextuelle par
    triangulation de trois méthodes complémentaires
    (Section~\ref{sec:trois-methodes}). \\
\textbf{3. Calcul ou bornage de $\mathsf{Conf}$, $\mathsf{Rel}$,
    $\mathsf{Opp}$} & Dériver les trois mesures à partir des propriétés
    cryptographiques publiées et de l'estimation de $H(C)$. \\
\textbf{4. Implications institutionnelles} & Interpréter le produit
    $\mathsf{Conf}\cdot\mathsf{Rel}\cdot\mathsf{Opp}$ et le situer
    dans un régime de capacité (Section~\ref{sec:regimes}). \\
\end{tableaucours}
\vspace{-0.8em}
\begin{center}{\small\itshape Tableau~28.1. Les quatre étapes de la
méthodologie CLO}\end{center}

\subsection{Les trois méthodes d'estimation de $H(C)$}
\label{sec:trois-methodes}

\begin{boxdef}
\textbf{Méthode~1. Bornes structurelles}~: compter les
caractéristiques contextuelles distinguables à partir des textes
légaux et réglementaires. \textit{Exemple~:} une juridiction
générale distingue $\approx 16$~types de preuves et $\approx 16$
contextes procéduraux, donnant
$H(C) \geq \log_2(16 \times 16) = 8$~bits.

\textbf{Méthode~2. Littérature d'expertise}~: s'appuyer sur les
résultats des sciences cognitives concernant la capacité de mémoire
de travail à long terme des experts (typiquement
$2^{12}$--$2^{15}$~``chunks'', cf.~Ericsson, cité
dans~\cite{minka2026clo}), suggérant une entropie contextuelle
d'expert de l'ordre de~$12$ à~$15$~bits.

\textbf{Méthode~3. Modélisation information-théorique}~: modéliser
la prise de décision institutionnelle comme une application
probabiliste de l'evidence vers les verdicts, et estimer les bornes
d'information mutuelle correspondantes.
\end{boxdef}

Ces trois méthodes convergent, selon~\cite{minka2026clo}, vers les
plages suivantes~: tribunaux généraux $H(C) \approx 8$--$16$~bits~;
régulateurs spécialisés $\approx 10$--$32$~bits~; experts de
domaine $\approx 20$--$64$~bits.

% ============================================================================
\section{Étude de Cas~1~: Helios. Vote Électronique}
\label{sec:cas-helios}
% ============================================================================

\subsection{Formalisation du protocole}

Helios~\cite{adida2009helios} utilise le chiffrement ElGamal et des
preuves à divulgation nulle de connaissance Fiat-Shamir sur des
groupes de logarithme discret~$2048$~bits.

\begin{itemize}
    \item \textbf{Messages}~: votes individuels $v \in \{1,\ldots,K\}$~;
    \item \textbf{Preuves}~: $\sigma = (\mathrm{Enc}(v), \pi_{\mathrm{valid}},
          \pi_{\mathrm{tally}})$, où $\mathrm{Enc}$ est le
          chiffrement ElGamal et $\pi$ les preuves Fiat-Shamir~;
    \item \textbf{Contexte}~: litige électoral devant un tribunal de
          juridiction générale, $H(C) \approx 8$~bits (Méthode~1~:
          $\approx 16$~types de preuves distinguables, $\approx
          16$~contextes procéduraux)~;
    \item \textbf{Verdicts}~: $\mathcal{V} = \{$élection valide,
          élection invalide, recomptage, demande d'investigation$\}$,
          soit $H(V \mid C) = 2$~bits.
\end{itemize}

\subsection{Calcul détaillé, étape par étape}

\begin{boxoutils}{Helios, calcul publié, reproduit intégralement}
\textbf{Confidentialité}. sécurité sémantique ElGamal sous
l'hypothèse DDH (\textit{Decisional Diffie-Hellman})~:
\[
\mathsf{Conf}(\text{Helios}) \approx 0{,}92
\]

\textbf{Fiabilité}. solidité des preuves à divulgation nulle,
erreur $< 2^{-80}$~:
\[
\mathsf{Rel}(\text{Helios}) \approx 0{,}95
\]

\textbf{Opposabilité}. décomposition explicite~:
les tribunaux peuvent extraire un \textbf{signal d'expert binaire}
depuis les preuves cryptographiques~: ``valide'' ou ``invalide''.
Cela représente
\[
I(V;\Sigma\mid C) \approx 1~\text{bit}
\]
sur les $H(V\mid C) = 2$~bits d'incertitude totale du verdict. D'où~:
\[
\mathsf{Opp}(\text{Helios}, C_{\text{tribunal}}) =
\frac{I(V;\Sigma\mid C)}{H(V\mid C)} \approx \frac{1}{2} = 0{,}50
\]

\textbf{Produit}~:
\[
\mathsf{Conf}\cdot\mathsf{Rel}\cdot\mathsf{Opp} \approx
0{,}92 \times 0{,}95 \times 0{,}50 \approx \mathbf{0{,}437}
\]
\end{boxoutils}

\textbf{Interprétation publiée}~: cette opposabilité de~$50\%$
signifie que les preuves cryptographiques résolvent la \emph{moitié}
de l'incertitude institutionnelle~; l'autre moitié requiert un
raisonnement juridique sur la conformité procédurale, l'intention
de l'électeur et le droit électoral spécifique à la
juridiction, des dimensions \textbf{inaccessibles à la seule
vérification cryptographique}.

% ============================================================================
\section{Étude de Cas~2~: zk-SNARKs. Conformité Financière}
\label{sec:cas-zksnark}
% ============================================================================

\subsection{Formalisation du protocole}

Référence~\cite{groth2016zksnark}~: preuves à divulgation nulle
succinctes (construction Groth16 sur courbe BN254) permettant à une
institution financière de prouver ``toutes les transactions
respectent les règles de conformité'' via une preuve de~200~octets,
sans révéler le détail des transactions.

\begin{itemize}
    \item \textbf{Contexte}~: audit réglementaire par une autorité
          financière, $H(C) \approx 11$~bits (expertise réglementaire~:
          $\approx 64$~catégories de preuves à travers
          $\approx 32$~cadres procéduraux,
          $\log_2(64\times 32) \approx 11$)~;
    \item \textbf{Verdicts}~: $\mathcal{V} = \{$conforme, non conforme,
          investigation requise, sanction, action corrective,
          escalade, approbation conditionnelle, audit forensique
          complet$\}$, soit $H(V\mid C) = \log_2(8) = 3$~bits.
\end{itemize}

\subsection{Calcul détaillé, étape par étape}

\begin{boxoutils}{zk-SNARK. calcul publié, reproduit intégralement}
\textbf{Confidentialité}. divulgation nulle calculatoire sous
hypothèses d'accouplement (\textit{pairing})~:
\[
\mathsf{Conf}(\text{zk-SNARK}) \approx 0{,}89
\]

\textbf{Fiabilité}. solidité de connaissance, en tenant compte
du \textit{trusted setup}~:
\[
\mathsf{Rel}(\text{zk-SNARK}) \approx 0{,}96
\]

\textbf{Opposabilité}. décomposition explicite en \textbf{deux}
bits d'information extractible~:
\begin{enumerate}
    \item \textit{Vérification calculatoire} ($1$~bit)~: la preuve
          vérifie ou non sous les contraintes d'accouplement~;
    \item \textit{Catégorisation de conformité} ($1$~bit)~:
          classification approximative visible dans la structure
          de la preuve (``passe toutes les contraintes'' vs.
          ``proche de la limite'' vs.\ ``signal de violation clair'').
\end{enumerate}
\[
I(V;\Sigma\mid C) \approx 2~\text{bits}, \qquad
\mathsf{Opp}(\text{zk-SNARK}, C_{\text{régulateur}}) =
\frac{2}{3} \approx 0{,}67
\]

\textbf{Produit}~:
\[
\mathsf{Conf}\cdot\mathsf{Rel}\cdot\mathsf{Opp} \approx
0{,}89 \times 0{,}96 \times 0{,}67 \approx \mathbf{0{,}572}
\]
\end{boxoutils}

\textbf{Interprétation publiée}~: l'opposabilité supérieure à celle
d'Helios ($67\%$ contre~$50\%$) reflète un \textbf{contexte
institutionnel plus riche} ($H(C)=11$~bits contre~$8$)~: les
preuves cryptographiques résolvent les deux tiers de l'incertitude
réglementaire. Mais le tiers restant, jugement réglementaire sur
la matérialité, évaluation du risque systémique, priorisation des
mesures d'exécution, demeure structurellement
\textbf{inaccessible}, ce qui constitue exactement la limite
documentée au Ch.~\ref{chap:for-cloud} de ce cours~: l'opacité
cryptographique qui protège la confidentialité des clients
empêche l'audit ciblé lorsqu'un schéma suspect émerge.

% ============================================================================
\section{Étude de Cas~3~: ECDSA. Signature Juridique}
\label{sec:cas-ecdsa}
% ============================================================================

\subsection{Formalisation du protocole}

Référence~\cite{johnson2001ecdsa}~: authentification standard via
ECDSA sur la courbe NIST~P-256 ($256$~bits de sécurité),
équation de vérification $sG = R + eP_{\mathrm{pub}}$.

\begin{itemize}
    \item \textbf{Contexte}~: litige contractuel civil avec
          infrastructure PKI établie, $H(C) \approx 20$~bits, décomposé
          en~:
          compréhension de l'infrastructure PKI
          ($\approx 2^{10}=1024$~types de certificats~: $10$~bits)~;
          précédents jurisprudentiels
          ($\approx 2^8=256$~catégories de précédents~: $8$~bits)~;
          cadres statutaires (ESIGN Act, eIDAS, réglementations
          étatiques~: $\approx 2^2=4$~cadres~: $2$~bits)~;
    \item \textbf{Verdicts}~: espace de décision riche couvrant
          validité, force d'authentification, poids juridique,
          recevabilité procédurale, validité temporelle, confirmation
          d'autorité, pertinence contextuelle, $H(V\mid C)
          \approx 20$~bits.
\end{itemize}

\subsection{Calcul détaillé, étape par étape}

\begin{boxoutils}{ECDSA, calcul publié, reproduit intégralement}
\textbf{Confidentialité}. les messages signés sont publics~;
seules les clés privées restent cachées~:
\[
\mathsf{Conf}(\text{ECDSA}) \approx 0{,}31
\]

\textbf{Fiabilité}. vérification déterministe, la falsification
exige de résoudre le problème du logarithme discret elliptique~:
\[
\mathsf{Rel}(\text{ECDSA}) \approx 0{,}98
\]

\textbf{Opposabilité}. décomposition en \textbf{cinq} dimensions
d'information extractible (contrairement aux zk-SNARKs qui
minimisent intentionnellement la divulgation, ECDSA permet aux
tribunaux d'évaluer indépendamment plusieurs dimensions
institutionnelles via les cadres juridiques et techniques établis)~:

\begin{tableaucours}{p{4.0cm}p{1.5cm}p{5.0cm}}{%
    \tableauHeader{Dimension extraite} &
    \tableauHeader{Bits} &
    \tableauHeader{Contenu}
}
Vérification de propriété de
    clé publique & 3 & Chaîne de certificats, confiance émetteur,
    restrictions d'usage \\
Liaison message-signature & 4 & Intégrité du hash, conformité
    de format, non-répudiation \\
Validité temporelle & 3 & Horodatage, période de validité du
    certificat, statut de révocation (CRL/OCSP) \\
Évaluation d'autorité légale & 3 & Vérification du signataire,
    formation du contrat, conformité juridictionnelle \\
Pertinence contextuelle & 5 & Type de transaction, conformité
    réglementaire, profil de risque \\
\textbf{Total} & \textbf{18} & \\
\end{tableaucours}
\vspace{-0.8em}
\begin{center}{\small\itshape Tableau~28.2. Décomposition de
$I(V;\Sigma\mid C)$ pour ECDSA (5~dimensions, total~18~bits)}\end{center}

\[
I(V;\Sigma\mid C) \approx 18~\text{bits}, \qquad
\mathsf{Opp}(\text{ECDSA}, C_{\text{tribunal civil}}) =
\frac{18}{20} \approx 0{,}90
\]

\textbf{Produit}~:
\[
\mathsf{Conf}\cdot\mathsf{Rel}\cdot\mathsf{Opp} \approx
0{,}31 \times 0{,}98 \times 0{,}90 \approx \mathbf{0{,}273}
\]
\end{boxoutils}

\textbf{Interprétation publiée}~: ECDSA atteint la plus haute
opposabilité des trois protocoles ($90\%$) précisément
\textbf{parce qu'il sacrifie la confidentialité par construction}~:
sa transparence permet aux tribunaux d'évaluer indépendamment
plusieurs dimensions à travers des cadres juridiques et techniques
établis de longue date (ESIGN Act, eIDAS, Ch.~11~(UE~SEC4031)
de ce cours). C'est, des trois protocoles, celui dont le produit
$\mathsf{Conf}\cdot\mathsf{Rel}\cdot\mathsf{Opp}$ est le \textbf{plus
faible} ($0{,}273$), illustrant que l'opposabilité maximale ne
garantit pas le meilleur compromis global, exactement la tension
documentée au Ch.~\ref{chap:CRO}.

% ============================================================================
\section{Tableau de Synthèse Publié}
\label{sec:tableau-synthese}
% ============================================================================

\begin{tableaucours}{p{3.0cm}p{1.6cm}p{1.6cm}p{1.6cm}p{2.4cm}}{%
    \tableauHeader{Système} &
    \tableauHeader{$\mathsf{Conf}$} &
    \tableauHeader{$\mathsf{Rel}$} &
    \tableauHeader{$\mathsf{Opp}$} &
    \tableauHeader{$\mathsf{Conf}\cdot\mathsf{Rel}\cdot\mathsf{Opp}$}
}
Helios (vote) & $0{,}92$ & $0{,}95$ & $0{,}50$ & $0{,}437$ \\
zk-SNARK (conformité) & $0{,}89$ & $0{,}96$ & $0{,}67$ & $0{,}572$ \\
ECDSA (signature) & $0{,}31$ & $0{,}98$ & $0{,}90$ & $0{,}273$ \\
\end{tableaucours}
\vspace{-0.8em}
\begin{center}{\small\itshape Tableau~28.3. Synthèse intégrale des trois
analyses publiées dans~\cite{minka2026clo}, Table~1}\end{center}

\begin{boiteRemarque}{Ce que ce tableau ne contient pas, et pourquoi}
Remarquez que ce tableau ne comporte que \textbf{trois} primitives~:
ce sont les seules pour lesquelles un calcul complet, avec
décomposition bit par bit publiée et vérifiable, existe dans la
littérature à la date de rédaction de ce cours. AES, ChaCha20,
RSA, Kyber, Dilithium, zk-STARKs et les signatures à seuil sont des
primitives importantes (plusieurs rencontrées au Ch.~\ref{chap:quantique}
de ce cours) mais \textbf{n'ont pas fait l'objet d'un calcul CLO
publié}. Leur attribuer des scores inventés, comme le faisait une
version antérieure de ce chapitre, constituerait exactement le type
de fausse précision scientifique dénoncé au Ch.~\ref{chap:CRO}.
La Section~\ref{sec:discussion-qualitative} les aborde, mais
\textbf{uniquement qualitativement}, en signalant explicitement
l'absence de calcul publié.
\end{boiteRemarque}

% ============================================================================
\section{Régimes de Capacité Institutionnelle}
\label{sec:regimes}
% ============================================================================

\subsection{Définition des trois régimes}

Le second article~\cite{minka2026conditional} introduit une
classification opérationnelle directement utile pour la conception
de systèmes~: pour un paramètre de sécurité $\lambda$ donné,
trois régimes de capacité institutionnelle $\Delta$ sont définis
selon la borne $\Delta/\lambda$.

\begin{tableaucours}{p{2.8cm}p{2.5cm}p{2.5cm}p{3.5cm}}{%
    \tableauHeader{Régime} &
    \tableauHeader{Plage de $\Delta$ ($\lambda=128$)} &
    \tableauHeader{Borne $\Delta/\lambda$} &
    \tableauHeader{Protocoles adaptés}
}
\textbf{Faible capacité} & $<13$~bits & $<0{,}10$ &
    Privilégier la transparence (style ECDSA)~; tribunaux de
    petites créances \\
\textbf{Capacité moyenne} & $13$--$64$~bits & $0{,}10$--$0{,}50$ &
    Protocoles équilibrés (style Helios)~; agences réglementaires
    avec personnel technique \\
\textbf{Capacité élevée} & $>64$~bits & $>0{,}50$ &
    Protocoles préservant la confidentialité (style zk-SNARK)~;
    cours de sécurité nationale, cellules de renseignement
    financier \\
\end{tableaucours}
\vspace{-0.8em}
\begin{center}{\small\itshape Tableau~28.4. Régimes de capacité
institutionnelle (d'après~\cite{minka2026conditional}, Table~1,
pour $\lambda=128$)}\end{center}

\subsection{Capacités $\Delta$ implicites~: validation croisée}

Une contribution remarquable du second article est de
\textbf{retrouver}, à partir des produits CLO publiés dans le
premier article, la capacité institutionnelle $\Delta$ implicite
de chaque contexte, une forme de validation croisée entre les
deux publications.

\begin{tableaucours}{p{2.8cm}p{2.2cm}p{1.5cm}p{2.5cm}p{2.0cm}}{%
    \tableauHeader{Système} &
    \tableauHeader{Produit CLO} &
    \tableauHeader{$\lambda$} &
    \tableauHeader{$\Delta$ implicite} &
    \tableauHeader{Régime}
}
Helios & $0{,}437$ & $128$ & $\gtrsim 56$~bits & Moyen-Élevé \\
zk-SNARK & $0{,}572$ & $128$ & $\gtrsim 73$~bits & Élevé \\
ECDSA & $0{,}273$ & $128$ & $\gtrsim 35$~bits & Moyen \\
\end{tableaucours}
\vspace{-0.8em}
\begin{center}{\small\itshape Tableau~28.5. Capacités $\Delta$
implicites, retrouvées à partir des produits CLO publiés
(d'après~\cite{minka2026conditional}, Table~2)}\end{center}

\begin{boiteRemarque}{La cohérence comme argument de validité}
Ces trois valeurs de $\Delta$ implicite (35, 56, 73~bits) tombent
\textbf{toutes} dans la plage cognitivement plausible $[20,80]$~bits
dérivée indépendamment à partir des travaux de Miller sur la mémoire
de travail (Ch.~\ref{chap:CRO}, Section~4). Plus précisément, les
auteurs notent un \textbf{facteur d'amplification} de~$1{,}8$
à~$3{,}7\times$ par rapport à la capacité individuelle de
base ($\approx 20$~bits), cohérent avec les facteurs de
documentation ($1{,}5\times$), de délibération collective
($1{,}5\times$) et de formation spécialisée ($1{,}3\times$)
identifiés au chapitre précédent. \textbf{Cette convergence entre
deux articles indépendants, utilisant des méthodes de calcul
différentes (calcul direct vs.\ inversion de la borne théorique),
constitue un argument de cohérence interne pour le cadre dans son
ensemble}, exactement le type de convergence multi-sources
valorisé depuis le Ch.~\ref{chap:for-sys} de ce cours.
\end{boiteRemarque}

% ============================================================================
\section{Discussion Qualitative~: Primitives Non Encore Analysées}
\label{sec:discussion-qualitative}
% ============================================================================

Cette section aborde, \textbf{uniquement de façon qualitative} et
en s'appuyant sur les propriétés cryptographiques établies présentées
au Ch.~\ref{chap:quantique}, des primitives rencontrées ailleurs
dans ce cours mais ne disposant pas d'un calcul CLO publié.

\begin{tableaucours}{p{2.8cm}p{4.0cm}p{4.5cm}}{%
    \tableauHeader{Primitive} &
    \tableauHeader{Positionnement qualitatif attendu} &
    \tableauHeader{Justification (sans calcul publié)}
}
AES-256 & Probablement proche de Helios en structure
    (forte $\mathsf{Conf}$/$\mathsf{Rel}$, $\mathsf{Opp}$ limitée) &
    Sans la clé, aucune vérification institutionnelle n'est
    possible~: structurellement analogue à la confidentialité
    ElGamal d'Helios \\
CRYSTALS-Kyber & Probablement comparable au zk-SNARK
    (haute $\mathsf{Conf}$, $\mathsf{Opp}$ contrainte par la
    nouveauté) & Standard récent (FIPS~203, Ch.~\ref{chap:quantique})~:
    jurisprudence et formation des décideurs encore limitées,
    plausiblement analogue à la situation des zk-SNARKs \\
zk-STARKs & Possiblement entre Helios et zk-SNARK & Transparent
    (pas de \textit{trusted setup}, contrairement aux zk-SNARKs)~:
    pourrait offrir une $\mathsf{Rel}$ légèrement supérieure, mais
    aucune donnée publiée ne permet de quantifier cet écart \\
\end{tableaucours}
\vspace{-0.8em}
\begin{center}{\small\itshape Tableau~28.6. Positionnement qualitatif,
sans valeurs numériques inventées}\end{center}

\begin{boxalert}
\textbf{Ce tableau est une invitation à la recherche, pas un résultat}

Les positionnements du Tableau~28.6 sont des \textbf{hypothèses
qualitatives raisonnables}, pas des résultats établis. Les calculer
rigoureusement (formalisation complète, estimation triangulée de
$H(C)$, décomposition bit par bit de $I(V;\Sigma\mid C)$) constitue
précisément le type de contribution que l'Exercice~28.3 de ce
chapitre vous invite à esquisser.
\end{boxalert}

% ============================================================================
\section{Implications pour la Conception de Systèmes Forensiques}
\label{sec:implications-conception-prim}
% ============================================================================

\subsection{Quatre étapes de sélection de protocole}

D'après~\cite{minka2026conditional}, Section~VI-C~:

\begin{enumerate}
    \item \textbf{Évaluer la capacité institutionnelle}~: estimer
          $\Delta$ via évaluation cognitive, disponibilité d'experts,
          complexité procédurale (méthode du Tableau~28.1, appliquée
          à l'institution plutôt qu'au protocole).
    \item \textbf{Identifier le régime cible}~: déterminer si le
          contexte relève du régime Faible, Moyen ou Élevé
          (Tableau~28.4).
    \item \textbf{Sélectionner le protocole adapté}~: régime
          faible~$\to$ privilégier la transparence (style ECDSA)~;
          régime moyen~$\to$ protocoles équilibrés (style Helios)~;
          régime élevé~$\to$ protocoles préservant la confidentialité
          (style zk-SNARK).
    \item \textbf{Renforcer la capacité si nécessaire}~: si le
          protocole souhaité excède la capacité institutionnelle
          disponible, investir dans la formation d'experts, la
          documentation technique, les processus de délibération
          collaborative, ou les mécanismes de témoignage expert
          externe (Ch.~12~(UE~SEC4031), experts agréés
          Décret~69/DF/544).
\end{enumerate}

\subsection{Application au contexte forensique camerounais}

\begin{boxscenario}{Positionner le système judiciaire camerounais dans
les régimes de capacité}
En s'appuyant sur les estimations de $H(C)$ du Tableau~25.2 du
Ch.~\ref{chap:bench} (juridiction générale~$\approx 8$~bits,
régulateurs spécialisés~$\approx 11$~bits) et sur le paramètre de
sécurité usuel $\lambda = 128$~bits~:

\begin{itemize}
    \item Un \textbf{tribunal de première instance camerounais
          non spécialisé} se situe, par analogie avec le calcul
          Helios ($H(C)\approx 8$~bits), dans un \textbf{régime
          de capacité faible à moyenne}~: $\Delta/\lambda \approx
          8/128 \approx 0{,}06$, proche de la borne du régime
          ``Faible''. \textbf{Implication}~: privilégier les
          preuves transparentes (signatures ECDSA, eIDAS,
          Ch.~11~(UE~SEC4031)) plutôt que des protocoles
          cryptographiquement sophistiqués mais opaques.
    \item Une \textbf{cellule spécialisée de lutte contre la
          cybercriminalité} (formée, avec accès à une expertise
          forensique dédiée, Ch.~12~(UE~SEC4031)), par analogie
          avec le contexte régulateur du zk-SNARK
          ($H(C)\approx 11$~bits), se rapproche du \textbf{régime
          moyen}~: $\Delta/\lambda \approx 11/128 \approx 0{,}09$.
    \item \textbf{Recommandation directement actionnable}~: la
          stratégie de formation des magistrats évoquée au
          Ch.~\ref{chap:bench} (Recommandation~3) n'est pas
          seulement souhaitable en abstrait, elle constitue,
          au sens précis du Tableau~28.4, le levier permettant de
          \textbf{faire progresser le système judiciaire camerounais
          d'un régime de capacité faible vers un régime moyen},
          élargissant ainsi l'éventail des protocoles cryptographiques
          dont les preuves resteraient opposables devant les
          tribunaux.
\end{itemize}
\end{boxscenario}

\separateur

\begin{boxresume}
\textbf{Ce qu'il faut retenir de ce chapitre~:}
\begin{itemize}
    \item La méthodologie CLO procède en \textbf{quatre étapes}~:
          formalisation, estimation triangulée de $H(C)$ (trois
          méthodes convergentes), calcul des trois mesures,
          interprétation institutionnelle.

    \item \textbf{Trois calculs intégralement publiés et reproduits}~:
          Helios ($0{,}92 \times 0{,}95 \times 0{,}50 = 0{,}437$)~;
          zk-SNARK ($0{,}89 \times 0{,}96 \times 0{,}67 = 0{,}572$)~;
          ECDSA ($0{,}31 \times 0{,}98 \times 0{,}90 = 0{,}273$).
          Chaque mesure d'opposabilité est décomposée bit par bit,
          dimension par dimension.

    \item \textbf{ECDSA atteint la plus haute opposabilité} ($90\%$)
          mais le \textbf{plus faible produit global} ($0{,}273$)~:
          la transparence maximale n'est pas le meilleur compromis
          général.

    \item \textbf{Trois régimes de capacité institutionnelle}
          (Faible~$<13$~bits, Moyen~$13$--$64$, Élevé~$>64$, pour
          $\lambda=128$) guident le choix de protocole~: transparence
          en régime faible, équilibre en régime moyen, confidentialité
          en régime élevé.

    \item \textbf{Validation croisée remarquable}~: les capacités
          $\Delta$ implicites retrouvées à partir des produits CLO
          publiés (35 à~73~bits) convergent avec l'estimation
          cognitive indépendante de Miller, cohérence interne
          entre les deux articles.

    \item \textbf{Rigueur méthodologique non négociable}~: seules
          les primitives effectivement calculées dans la littérature
          publiée reçoivent des valeurs numériques dans ce chapitre.
          Les autres (AES, Kyber, Dilithium, zk-STARKs) sont
          discutées qualitativement, sans invention de scores.

    \item \textbf{Application camerounaise}~: la formation des
          magistrats (Ch.~\ref{chap:bench}) constitue, dans le
          vocabulaire précis de ce chapitre, le levier permettant
          de faire progresser le système judiciaire d'un régime
          de capacité faible vers un régime moyen.
\end{itemize}
\end{boxresume}

\bigskip

\begin{boxexo}[title={\ding{45}\quad Exercice~28.1. Reproduire un calcul publié \niveaubase}]
\begin{enumerate}[(a)]
    \item Pour le cas zk-SNARK (Section~\ref{sec:cas-zksnark}),
          vérifiez par le calcul que
          $0{,}89 \times 0{,}96 \times 0{,}67 \approx 0{,}572$.
    \item Si l'estimation de $H(C)$ pour le régulateur passait
          de~$11$ à~$16$~bits (régulateur encore plus spécialisé),
          et si $I(V;\Sigma\mid C)$ restait à~$2$~bits, quelle
          serait la nouvelle valeur de $\mathsf{Opp}$~? Le produit
          global augmenterait-il ou diminuerait-il~? Justifiez par
          le calcul.
    \item Expliquez, en une phrase, pourquoi $H(V\mid C)$ pour
          ECDSA ($20$~bits) est nettement supérieur à celui d'Helios
          ($2$~bits), en vous appuyant sur la nature des espaces de
          verdicts respectifs.
\end{enumerate}
\end{boxexo}

\begin{boxexo}[title={\ding{45}\quad Exercice~28.2. Régimes de capacité appliqués \niveaumoyen}]
\begin{enumerate}[(a)]
    \item Pour $\lambda = 256$~bits (sécurité renforcée, pertinente
          face aux préoccupations post-quantiques du
          Ch.~\ref{chap:quantique}), recalculez les bornes
          $\Delta/\lambda$ des trois régimes du Tableau~28.4.
          Qu'observez-vous sur la difficulté relative à atteindre
          le régime ``Élevé'' à mesure que $\lambda$ augmente~?
    \item Un laboratoire forensique national nouvellement créé
          dispose d'experts formés mais d'une jurisprudence encore
          limitée. Dans quel régime le situeriez-vous~? Quel type
          de protocole recommanderiez-vous en priorité~?
    \item En vous appuyant sur le Tableau~28.5, expliquez pourquoi
          la convergence entre les capacités $\Delta$ implicites
          et l'estimation cognitive de Miller constitue un argument
          scientifique pertinent, et quelles réserves méthodologiques
          devraient néanmoins accompagner cette interprétation
          (référez-vous à l'avertissement méthodologique de ce
          chapitre).
\end{enumerate}
\end{boxexo}

\begin{boxexo}[title={\ding{45}\quad Exercice~28.3. Esquisser un calcul original \niveauexpert}]
Choisissez une primitive \textbf{non couverte} par les publications
analysées dans ce chapitre~: AES-256, CRYSTALS-Kyber, ou
CRYSTALS-Dilithium (toutes rencontrées au Ch.~\ref{chap:quantique}),
appliquée à un contexte forensique camerounais de votre choix
(par exemple~: chiffrement BitLocker d'un disque saisi, signature
post-quantique d'un rapport d'expertise).
\begin{enumerate}[(a)]
    \item Formalisez le protocole selon les quatre composantes
          $\mathcal{M}, \Sigma, \mathcal{C}, \mathcal{V}$
          (Étape~1 de la méthodologie, Section~\ref{sec:methodologie}).
    \item Estimez $H(C)$ pour votre contexte choisi en appliquant
          au moins deux des trois méthodes
          (Section~\ref{sec:trois-methodes}), en justifiant vos
          hypothèses aussi explicitement que les calculs publiés
          analysés dans ce chapitre.
    \item Proposez une décomposition dimension par dimension de
          $I(V;\Sigma\mid C)$, sur le modèle du Tableau~28.2
          (ECDSA), avec au moins trois dimensions justifiées.
    \item Calculez $\mathsf{Opp}$ et situez votre primitive dans
          un régime de capacité (Tableau~28.4). Comparez votre
          résultat aux trois protocoles publiés
          (Tableau~28.3), votre estimation est-elle cohérente
          avec le positionnement qualitatif suggéré au
          Tableau~28.6~?
    \item Rédigez un paragraphe de conclusion identifiant
          explicitement les limites de votre propre estimation
          (sur le modèle de l'avertissement méthodologique des
          auteurs en ouverture de ce chapitre), qu'est-ce qui
          rendrait votre calcul rigoureux au sens d'une publication
          évaluée par les pairs, et qu'est-ce qui en l'état reste
          une approximation pédagogique~?
\end{enumerate}
\end{boxexo}

\bigskip

\begin{boxinfo}
\textbf{Transition.} Ce chapitre clôt l'exploration théorique du
Trilemme CRO entamée au Chapitre~\ref{chap:CRO}. Vous disposez
désormais des outils mathématiques (théorie de l'information), des
résultats établis (théorèmes prouvés et publiés) et de la méthode
de calcul reproductible nécessaires pour appliquer ce cadre à
n'importe quel système de preuve numérique que vous rencontrerez
dans votre pratique professionnelle, qu'il s'agisse d'un système
existant ou d'un protocole que vous serez amené à concevoir ou
évaluer.
\end{boxinfo}

% ============================================================================
% Références bibliographiques propres à ce chapitre
% (entrées supplémentaires à intégrer dans D_bibliographie.tex,
% en complément de celles du Ch.26)
% ============================================================================
% @inproceedings{groth2016zksnark,
%   author    = {Groth, Jens},
%   title     = {On the Size of Pairing-Based Non-Interactive Arguments},
%   booktitle = {Proc. Annu. Int. Conf. Theory Appl. Cryptograph. Techn.},
%   pages     = {305--326},
%   year      = {2016}
% }